\subsection{Actual native images and lifts of integral scheme morphisms}
\label{subsec:native-integral-images}

\begin{theorem}[Genuine integral-image lifting of whole native points]
\label{scheme:integral-native-image}
\lean{EnvelopeProperty.Berkovich.SchemeValuedPoint.exists_map_eq_of_point_mem_integralImage,
EnvelopeProperty.Berkovich.SchemeValuedPoint.range_map_of_isIntegralHom,
EnvelopeProperty.Berkovich.SchemeValuedPoint.map_surjective_of_isIntegralHom}
\leanok
\uses{spectrum:integral-image,scheme:native-closed-immersion,
scheme:affine-chart-section-values}
For an actual integral scheme morphism $f:X\to Y$, a native valued point
$y$ of $Y$ has a whole native preimage precisely when its actual scheme
support lies in the genuine scheme image of $f$. Thus
\[
 \operatorname{im}(f^{\mathrm{val}})
 =\{y:\operatorname{supp}(y)\in\operatorname{im}(f)\}.
\]
If the actual scheme-point map is surjective, so is its native map.
No field, coefficient bound, finite-type, noetherian, separatedness or
completeness hypothesis is required.
\end{theorem}
\begin{proof}
Necessity is the original definition of the native support map. For
sufficiency, let $x\in X$ be an actual scheme preimage witness for the
support of $y$, and choose a genuine affine target open $U$ containing
that support. Its preimage under the integral morphism is affine,
and the original section homomorphism
$\alpha:\Gamma(Y,U)\to\Gamma(X,f^{-1}U)$ is integral.

The native affine chart of $y$ is a real seminorm $z$ of the target
section ring. If an actual section $a$ lies in $\ker\alpha$, then its
source residue evaluation at $x$ is zero. Genuine scheme evaluation
naturality identifies this zero with the target residue evaluation of
$a$ pulled back under the actual residue-field embedding
$\kappa(f(x))\to\kappa(x)$. That embedding is injective, so the original
target residue element is zero. The original valuation of $y$ therefore
vanishes on $a$. This proves the actual kernel obstruction
$\ker\alpha\subseteq\operatorname{supp}(z)$ using equality of residue
elements, rather than only numerical generator values.

The independently constructed integral-spectrum image theorem produces
an actual whole seminorm of the source section ring pulling back to
$z$. Take its inverse genuine source native affine chart. The entire
native affine chart square identifies the native image's whole target
section seminorm with $z$. Injectivity of the actual target chart then
recovers the original whole native point $y$. The chosen scheme point
$x$ was only used to prove the obstruction; no assertion is made that
the constructed valuation has that particular source support.

This proves sufficiency and the exact image equality. Actual
scheme-point surjectivity supplies a support preimage witness for every
native point, proving native surjectivity without a native lift
certificate or additional point-extension axiom.
\end{proof}

\begin{theorem}[Exact integral images above a common prescribed base point]
\label{scheme:relative-integral-native-image}
\lean{EnvelopeProperty.Berkovich.RelativeSchemeValuedPoint.exists_map_eq_of_point_mem_integralImage,
EnvelopeProperty.Berkovich.RelativeSchemeValuedPoint.range_map_of_isIntegralHom,
EnvelopeProperty.Berkovich.RelativeSchemeValuedPoint.map_surjective_of_isIntegralHom}
\leanok
\uses{scheme:integral-native-image,scheme:relative-native-closed-immersion}
For an actual commuting integral morphism over $\operatorname{Spec}R$
and any prescribed real seminorm $v$ on $R$, a native point above $v$
has a native preimage above that same whole base point exactly when its
support lies in the actual scheme image. Scheme-point surjectivity
therefore implies native fixed-base surjectivity. No bound on $v$ or
proper or separated structure morphism is needed.
\end{theorem}
\begin{proof}
First construct the whole absolute lift by the preceding theorem.
The real structure diagram and whole native composition identity
identify the source fixed-base locus with the target-locus preimage.
Consequently the actual lift lies above the same entire prescribed
base point $v$. Restrict it to that source subtype; equality of its
original entire native image gives the required whole relative point
equality. Necessity and actual scheme surjectivity then give the two
remaining conclusions.

This is a common-coefficient theorem. It does not claim surjectivity
for an arbitrary differently prescribed source coefficient seminorm
under a general coefficient change, nor construct a normalization,
curve or analytic potential.
\end{proof}
